\chapter{General Introduction}
\begin{spacing}{1.3}

Machine learning has undergone a profound transformation over the past decade. What began as a domain of academic research, confined to notebooks and controlled laboratory environments, has matured into a critical component of enterprise software systems. The narrative has shifted: organisations no longer ask whether they can build a model, but whether they can deploy, monitor, and maintain it at scale. This evolution mirrors the broader trajectory of software engineering, where the introduction of DevOps practices revolutionised the reliability and velocity of software delivery.

Despite significant advances in model architectures, frameworks, and hardware accelerators, the bridge between a trained model and a production system remains fragile. A frequently cited statistic in the MLOps community is that 87\% of machine learning models never reach production. The reasons are manifold: brittle data pipelines, unreproducible training procedures, absence of monitoring, model decay over time, and a cultural divide between data scientists and software engineers. Sculley et al.~\cite{sculley2015} described this phenomenon as \textit{hidden technical debt} in machine learning systems, arguing that the ML component itself is a tiny fraction of the overall system, surrounded by vast infrastructure for data collection, feature extraction, configuration, monitoring, and serving.

MLOps -- a portmanteau of Machine Learning and Operations -- extends the principles of DevOps to the unique challenges of ML systems. At its core, MLOps advocates for automated testing and validation of data, code, and models on every change (Continuous Integration), automated deployment pipelines moving models from training to serving (Continuous Delivery), automated retraining triggered by data drift (Continuous Training), versioning of data, models, experiments, and configurations, and production observability of model performance and system health. The MLOps maturity model proposed by Google~\cite{google_mlops} defines three levels, from manual script-driven processes to full CI/CD automation.

Healthcare applications occupy a uniquely demanding position in the machine learning landscape. Unlike recommendation systems where failure results in lost revenue, failures in medical AI can result in misdiagnosis and patient harm. The stakes demand safety mechanisms that degrade gracefully under uncertainty, full reproducibility where every prediction is traceable to its training pipeline, and continuous monitoring to detect silent model degradation. MLOps provides the operational framework to satisfy these requirements, making it not a luxury but a prerequisite in medical AI.

Melanoma is the deadliest form of skin cancer, accounting for approximately 75\% of skin cancer-related deaths despite representing only 1\% of diagnoses. The prognosis is starkly bifurcated by stage at detection: early-stage detection yields a 5-year survival rate exceeding 99\%, while late-stage detection drops survival below 30\%. Dermoscopy improves diagnostic accuracy over naked-eye examination but requires significant training. Dermatologists are in short supply globally, with particularly acute shortages in rural and underserved communities. Automated screening systems deployed on edge devices -- smartphones, tablets, or Raspberry Pi devices -- offer a compelling solution by enabling offline, low-cost, real-time screening without cloud dependency, preserving data privacy in the process.

This project is a research and educational prototype, not a medical device. It has not been evaluated by any regulatory authority. The pipeline is designed with safety guardrails -- uncertainty quantification, confidence thresholds, and mandatory clinical review flags -- that align with regulatory best practices. The HAM10000 dataset~\cite{tschandl2018}, comprising 10,015 dermoscopic images across seven diagnostic categories, serves as the training and evaluation foundation.

The present report documents the design, implementation, and evaluation of a complete 12-phase MLOps pipeline for melanoma detection, spanning data ingestion through edge deployment. Chapter~2 presents the clinical context, dataset characteristics, pipeline architecture, and model architectures. Chapter~3 reviews the state of the art in MLOps platforms, melanoma detection literature, and edge inference optimisation. Chapter~4 defines functional, non-functional, and MLOps-specific requirements. Chapter~5 provides a detailed, phase-by-phase account of the entire implementation. Chapter~6 presents model performance results and the Gradio frontend interface. Chapter~7 assesses project achievements against requirements, and the conclusion outlines directions for future work.

\end{spacing}
